\chapter[Representabilitat i Yoneda]{\texorpdfstring{Propietats universals,\\ representabilitat i Yoneda}{Propietats universals, representabilitat i Yoneda}}

Aquest test recull els conceptes essencials del capítol. Cada pregunta té una única resposta correcta.

\begin{Form}
\iniciTest{yon}{b,c,b,c,b,d,c,b,c,b,d,c,b,c}

\begin{enumerate}

\pregunta Quan és \textbf{terminal} un objecte $1$ d'una categoria?
\begin{enumerate}
  \opcio{a}{Quan hi ha exactament un morfisme $1\to X$ per a cada $X$.}
  \opcio{b}{Quan hi ha exactament un morfisme $X\to 1$ per a cada $X$.}
  \opcio{c}{Quan no hi ha cap morfisme cap a ell.}
  \opcio{d}{Quan és isomorf a tots els altres objectes.}
\end{enumerate}
\feedback

\pregunta Per què l'objecte terminal és únic llevat d'isomorfisme?
\begin{enumerate}
  \opcio{a}{Perquè només n'hi pot haver un literalment.}
  \opcio{b}{Perquè tots els objectes són terminals.}
  \opcio{c}{Perquè entre dos terminals els morfismes únics es componen donant les identitats, i per tant són isomorfismes.}
  \opcio{d}{Perquè l'objecte terminal no existeix mai.}
\end{enumerate}
\feedback

\pregunta Una \textbf{propietat universal} caracteritza un objecte...
\begin{enumerate}
  \opcio{a}{per la seva constitució interna.}
  \opcio{b}{per la manera com es relaciona amb tots els altres objectes, via un functor representable.}
  \opcio{c}{per la seva cardinalitat.}
  \opcio{d}{només a $\cat{Set}$.}
\end{enumerate}
\feedback

\pregunta Que $1$ sigui terminal equival a dir que:
\begin{enumerate}
  \opcio{a}{$\Hom(1,-)$ és constant.}
  \opcio{b}{$\Hom(-,1)$ no és functor.}
  \opcio{c}{$\Hom(-,1)$ és isomorf al functor constant amb valor un punt.}
  \opcio{d}{$\Hom(-,1)$ és buit.}
\end{enumerate}
\feedback

\pregunta Quan es diu que un functor $P\colon\cat{C}\op\to\cat{Set}$ és \textbf{representable}?
\begin{enumerate}
  \opcio{a}{Quan és constant.}
  \opcio{b}{Quan és isomorf a $\Hom(-,A)$ per a algun objecte $A$.}
  \opcio{c}{Quan és ple i fidel.}
  \opcio{d}{Quan té un sol element a cada conjunt.}
\end{enumerate}
\feedback

\pregunta Què és la \textbf{immersió de Yoneda}?
\begin{enumerate}
  \opcio{a}{Un objecte terminal.}
  \opcio{b}{Una transformació natural.}
  \opcio{c}{Un prefeix qualsevol.}
  \opcio{d}{El functor $\yon\colon\cat{C}\to\cat{Set}^{\cat{C}\op}$ amb $\yon(A)=\Hom(-,A)$.}
\end{enumerate}
\feedback

\pregunta Els objectes de $\cat{Set}^{\cat{C}\op}$ s'anomenen:
\begin{enumerate}
  \opcio{a}{functors oblit.}
  \opcio{b}{cons.}
  \opcio{c}{prefeixos.}
  \opcio{d}{monoides.}
\end{enumerate}
\feedback

\pregunta Què afirma el \textbf{lema de Yoneda}?
\begin{enumerate}
  \opcio{a}{Que tota categoria és petita.}
  \opcio{b}{Que $\Nat(\Hom(-,A),P)\cong P(A)$, naturalment en $A$ i $P$.}
  \opcio{c}{Que tot functor és representable.}
  \opcio{d}{Que $\Hom(A,B)$ és sempre buit.}
\end{enumerate}
\feedback

\pregunta A la demostració del lema, a quina dada s'envia una transformació natural $\eta\colon\Hom(-,A)\Rightarrow P$?
\begin{enumerate}
  \opcio{a}{A $\eta_A$.}
  \opcio{b}{A tot el functor $P$.}
  \opcio{c}{A l'element $\eta_A(\id_A)\in P(A)$.}
  \opcio{d}{A la identitat de $P$.}
\end{enumerate}
\feedback

\pregunta Una conseqüència directa del lema de Yoneda és que la immersió $\yon$ és:
\begin{enumerate}
  \opcio{a}{constant.}
  \opcio{b}{plena i fidel.}
  \opcio{c}{essencialment exhaustiva.}
  \opcio{d}{un isomorfisme de categories.}
\end{enumerate}
\feedback

\pregunta Segons el corol·lari de Yoneda, $A\cong B$ si i només si:
\begin{enumerate}
  \opcio{a}{$A$ i $B$ tenen la mateixa cardinalitat.}
  \opcio{b}{$A=B$.}
  \opcio{c}{hi ha un functor entre ells.}
  \opcio{d}{$\Hom(-,A)\cong\Hom(-,B)$ com a prefeixos.}
\end{enumerate}
\feedback

\pregunta L'objecte que representa un functor representable és:
\begin{enumerate}
  \opcio{a}{sempre $\cat{Set}$.}
  \opcio{b}{mai únic.}
  \opcio{c}{únic llevat d'isomorfisme.}
  \opcio{d}{un objecte inicial.}
\end{enumerate}
\feedback

\pregunta La immersió de Yoneda, en general, \emph{no} és:
\begin{enumerate}
  \opcio{a}{fidel.}
  \opcio{b}{essencialment exhaustiva (no tot prefeix és representable).}
  \opcio{c}{plena.}
  \opcio{d}{un functor.}
\end{enumerate}
\feedback

\pregunta Quina és la lectura intuïtiva del corol·lari de Yoneda?
\begin{enumerate}
  \opcio{a}{Que els objectes no es poden distingir.}
  \opcio{b}{Que tot objecte és terminal.}
  \opcio{c}{Que un objecte queda determinat, llevat d'isomorfisme, per la xarxa de morfismes que hi arriben.}
  \opcio{d}{Que els morfismes no importen.}
\end{enumerate}
\feedback

\end{enumerate}
\clauestatica
\finalitzaTest
\end{Form}
